\begin{titlepage}
  \begin{center}

  {\Huge AXIS\_UART}

  \vspace{25mm}

  \includegraphics[width=0.90\textwidth,height=\textheight,keepaspectratio]{img/AFRL.png}

  \vspace{25mm}

  \today

  \vspace{15mm}

  {\Large Jay Convertino}

  \end{center}
\end{titlepage}

\tableofcontents

\newpage

\section{Usage}

\subsection{Introduction}

\par
Simple UART core for TTL rs232 software mode data communications. No hardware handshake.
This contains its own internal baud rate generator that creates an enable to allow data output
or sampling. Baud clock and aclk can be the same clock.

RTS/CTS is not implemented, it simply asserts it as if its always ready, and ignores CTS.

\subsection{Dependencies}

\par
The following are the dependencies of the cores.

\begin{itemize}
  \item fusesoc 2.X
  \item iverilog (simulation)
  \item cocotb (simulation)
\end{itemize}

\subsubsection{fusesoc\_info Depenecies}
\begin{itemize}
\item dep
	\begin{itemize}
	\item AFRL:clock:mod\_clock\_ena\_gen:1.0.1
	\item AFRL:utility:helper:1.0.0
	\end{itemize}
\item dep\_tb
	\begin{itemize}
	\item AFRL:simulation:axis\_stimulator
	\item AFRL:utility:sim\_helper
	\end{itemize}
\end{itemize}


\subsection{In a Project}
\par
This core connects a UART to the AXIS bus. Meaning this is a streaming device only. Connect the RX/TX to the UART in question and connect the AXIS to its intended endpoints.

\section{Architecture}
\par
This core is made up of other cores that are documented in detail in there source. The cores this is made up of are the,
\begin{itemize}
  \item \textbf{axis\_uart\_tx} Interface with UART TX and present the data over AXIS interface (see core for documentation).
  \item \textbf{axis\_uart\_rx} Interface with UART RX and present the data over AXIS interface (see core for documentation).
  \item \textbf{uart\_baud\_gen} Generates BAUD clock for RX and TX based on modulo divide method (see core for documentation).
\end{itemize}

\section{Building}

\par
The AXIS UART is written in Verilog 2001. It should synthesize in any modern FPGA software. The core comes as a fusesoc packaged core and can be
included in any other core. Be sure to make sure you have meet the dependencies listed in the previous section.

\subsection{fusesoc}
\par
Fusesoc is a system for building FPGA software without relying on the internal project management of the tool. Avoiding vendor lock in to Vivado or Quartus.
These cores, when included in a project, can be easily integrated and targets created based upon the end developer needs. The core by itself is not a part of
a system and should be integrated into a fusesoc based system. Simulations are setup to use fusesoc and are a part of its targets.

\subsection{Source Files}

\subsubsection{fusesoc\_info File List}
\begin{itemize}
\item src
	\begin{itemize}
	\item src/axis\_spi\_master.v
	\end{itemize}
\item tb
	\begin{itemize}
	\item tb/tb\_spi.v
	\end{itemize}
\item tb\_cocotb
	\begin{itemize}
	\item {'tb/tb\_cocotb.py': {'file\_type': 'user', 'copyto': '.'}}
	\item {'tb/tb\_cocotb.v': {'file\_type': 'verilogSource'}}
	\end{itemize}
\end{itemize}


\subsection{Targets}

\subsubsection{fusesoc\_info Targets}
\begin{itemize}
\item default
	\begin{itemize}
	\item[$\space$] Info: Default for IP intergration.
	\item src
	\item dep
	\end{itemize}
\item sim
	\begin{itemize}
	\item[$\space$] Info: Base simulation using icarus as default.
	\item src
	\item dep
	\item tb
	\item dep\_tb
	\item IN\_FILE\_NAME
	\item OUT\_FILE\_NAME
	\item RAND\_READY
	\end{itemize}
\item sim\_rand\_data
	\begin{itemize}
	\item[$\space$] Info: Use random data as sim input.
	\item src
	\item dep
	\item tb
	\item dep\_tb
	\item IN\_FILE\_NAME=random.bin
	\item OUT\_FILE\_NAME=out\_random.bin
	\item RAND\_READY
	\item FIFO\_DEPTH
	\end{itemize}
\item sim\_rand\_ready\_rand\_data
	\begin{itemize}
	\item[$\space$] Info: Use random data with a random ready as sim input.
	\item src
	\item dep
	\item tb
	\item dep\_tb
	\item IN\_FILE\_NAME=random.bin
	\item OUT\_FILE\_NAME=out\_random.bin
	\item RAND\_READY=1
	\item FIFO\_DEPTH
	\end{itemize}
\item sim\_8bit\_count\_data
	\begin{itemize}
	\item[$\space$] Info: Use counter data as sim input.
	\item src
	\item dep
	\item tb
	\item dep\_tb
	\item IN\_FILE\_NAME=8bit\_count.bin
	\item OUT\_FILE\_NAME=out\_8bit\_count.bin
	\item RAND\_READY
	\item FIFO\_DEPTH
	\end{itemize}
\item sim\_rand\_ready\_8bit\_count\_data
	\begin{itemize}
	\item[$\space$] Info: Use counter data with a random ready as sim input.
	\item src
	\item dep
	\item tb
	\item dep\_tb
	\item IN\_FILE\_NAME=8bit\_count.bin
	\item OUT\_FILE\_NAME=out\_8bit\_count.bin
	\item RAND\_READY=1
	\item FIFO\_DEPTH
	\end{itemize}
\item sim\_baud
	\begin{itemize}
	\item[$\space$] Info: Simulate only the baud rate generator.
	\item src
	\item dep
	\item tb
	\item dep\_tb
	\item IN\_FILE\_NAME
	\item OUT\_FILE\_NAME
	\item RAND\_READY
	\end{itemize}
\item sim\_rx
	\begin{itemize}
	\item[$\space$] Info: Simulate only the rx block.
	\item src
	\item dep
	\item tb
	\item dep\_tb
	\item IN\_FILE\_NAME
	\item OUT\_FILE\_NAME
	\item RAND\_READY
	\end{itemize}
\item sim\_tx
	\begin{itemize}
	\item[$\space$] Info: Simulate only the tx block.
	\item src
	\item dep
	\item tb
	\item dep\_tb
	\item IN\_FILE\_NAME
	\item OUT\_FILE\_NAME
	\item RAND\_READY
	\end{itemize}
\end{itemize}


\subsection{Directory Guide}

\par
Below highlights important folders from the root of the directory.

\begin{enumerate}
  \item \textbf{docs} Contains all documentation related to this project.
    \begin{itemize}
      \item \textbf{manual} Contains user manual and github page that are generated from the latex sources.
    \end{itemize}
  \item \textbf{src} Contains source files for the core
  \item \textbf{tb} Contains test bench files for iverilog and cocotb
    \begin{itemize}
      \item \textbf{cocotb} testbench files
    \end{itemize}
\end{enumerate}

\newpage

\section{Simulation}
\par
There are a few different simulations that can be run for this core.

\subsection{iverilog}
\par
iverilog is used for simple test benches for quick verification, visually, of the core.

\subsection{cocotb}
\par
Future simulations will use cocotb. This feature is not yet implemented.

\newpage

\section{Module Documentation} \label{Module Documentation}

\begin{itemize}
\item \textbf{axis\_uart\_tx} Interfaces AXIS to the UART transmit line.\\
\item \textbf{axis\_uart\_rx} Interfaces AXIS to the UART receive line.\\
\item \textbf{uart\_baud\_gen} Generates a Baud rate clock for the UART based on the input clock using modulo division method.
The modulo method works by loading the clock speed (in hertz) used for generaton into a register. This register is then
subtracted from by the desired baud clock (in hertz). Once the counter is depleted the remainder, if it exits, is added
to the counter plus the original clock speed and the operation repeats. Each remiander is added back to the counter. This
allows non-even divisions to evetually average out to the target rate.\\
\item \textbf{axis\_uart Wrapper} for all of the above modules to create a singular device to interface with.\\
\end{itemize}


